\chapter*{How to use this book}
\addcontentsline{toc}{chapter}{How to use this book}
\markboth{How to use this book}{}

\section*{Who it is for}

An experienced engineer who is not an AI Engineer by title and is aiming at
Staff level. You have shipped things, you have operated them, and the gap you
are closing is a domain gap and an evidence gap rather than an engineering one.

If you are early in your career, or if you want to train models rather than
build systems on top of them, this is the wrong book and
Chapter~\ref{ch:three-jobs} says why in more detail.

\section*{What is in it}

\begin{kittable}{@{}llX@{}}
\textbf{Part} & \textbf{Pages} & \textbf{Read it} \\
\midrule
I The map & short & first, and before you study anything \\
II The evidence & short & before you apply \\
III The process & medium & when a process starts \\
IV The build rounds & long & when a build round is scheduled \\
V The bank & long & in the order the chapters are in \\
\end{kittable}

Part~V is ordered by how often each surface is reported to be asked, not by
how interesting it is. That ordering is itself the advice: study it in the
order it is printed, and Chapter~\ref{ch:noise} says what to leave out.

Alongside the book: \code{templates/} to copy and fill in per company,
\code{onepagers/} to print before a call, and \code{drills/} for the rehearsal
in Chapter~\ref{ch:rehearsal}. Appendices~\ref{app:templates}
and~\ref{app:drills} index them.

\section*{How claims are marked}

Almost nothing in this subject can be measured, so the kind of every claim is
marked on the page.

\begin{itemize}
  \item \reported{A statement about what the market does carries a key in the
        margin of the sentence.}{field-guide} The key resolves in
        \code{sources.json}, which records the URL, the date it was checked,
        how good the source is, and whether it was read directly or through
        somebody else's summary.
  \item \reported[single]{A claim resting on a single source says so, in the
        sentence, like this.}{kalra-transition} The book does not let you find
        that out by opening a file.
  \item \judgement{A statement marked like this is the author's opinion. It is
        not sourced, because a citation on an opinion is borrowed authority,
        and you are invited to disagree with it.}
\end{itemize}

Every chapter ends with \emph{What this chapter does not claim}. It is not
modesty; it is the part that tells you how much weight the chapter will bear.

\section*{What this book refuses to do}

\textbf{No figures it cannot source.} No salary bands, no premiums, no pass
rates. Those numbers are regional, they are stale within two quarters, and the
versions in circulation are almost all unsourced. What is here instead is the
method for measuring your own market, and a template to record it in.

\textbf{No case study, and no company.} The method is general or it is not
worth having. A build gate refuses any named employer, any individual, and any
currency figure, so this cannot drift back.

\textbf{No claims it could not trace.} Appendix~\ref{app:refused} lists the
figures that every competing resource repeats, names each one, and says why it
is not here. \judgement{That appendix is the most useful thing in the book, and
if you read only one page of it, read that one.}

\section*{If you have one week}

\begin{kitmethod}
\textbf{Day 1.} Chapters~\ref{ch:three-jobs} to \ref{ch:segment}, then the fit
check in Chapter~\ref{ch:fit} for the role in front of you. If it comes back
red, stop \dash{} you have saved the week.

\textbf{Days 2--4.} Build one small thing end to end, deploy it, and put an
evaluation harness on it (Chapters~\ref{ch:portfolio} and \ref{ch:evals}).
Spend the last hour on the cost arithmetic, out loud.

\textbf{Day 5.} Part~V, chapters \ref{ch:retrieval} to \ref{ch:cost}, in order.

\textbf{Day 6.} The rehearsal in Chapter~\ref{ch:rehearsal}, full length, with
another person. This is the highest-value day and it is the one that gets
dropped, because it needs somebody else's diary.

\textbf{Day 7.} The claim inventory and the banned-claims list
(Chapter~\ref{ch:claims}), the interviewer briefs
(Chapter~\ref{ch:person}), and print the cards.
\end{kitmethod}

\section*{A note on its own evidence}

This book was assembled from a market survey with a stated method, published
engineering ladders, practitioner writing, and one author's experience of
running a process of this kind end to end.

That last source is not a case study and does not appear as one. It is why
certain chapters exist \dash{} the fit check, the evidence classes, the
banned-claims list \dash{} and it is n{\kern0.08em=\kern0.08em}1, so nothing in
this book rests on it. Where a chapter is the author's judgement, the chapter
says so.

\clearpage
